%% kapa: exam=1 round=22 year=2011 period=1 form=A subject=civil_law subject_ko=민법 seq=5 number=5
%% source: Q-Net 감정평가사 기출문제 (cst00309) artlSeq=5036790 file=extracted/1008/5036790/2011년 제22회 감정평가사 1차시험 기출문제/1교시.hwp
%% conversion: kordoc 4.24.51 -> tools/gen1.py (auto); review=proofread-sonnet; flags=none
\begin{question}{5}
복대리에 관한 설명으로 옳은 것은?
\begin{choices}
  \item 복대리인은 대리인이 본인의 이름으로 선임한 본인의 대리인이다.
  \item 대리인이 복대리인을 선임한 후에는 대리인의 대리권은 소멸한다.
  \item 법정대리인에 의해 선임된 복대리인은 법정대리인이다.
  \item 복대리인이 제3자와 계약을 체결하는 경우, 본인의 이름 이외에 대리인의 이름도 현명하여야 한다.
  \item 복대리인은 대리인의 대리권에 기초하기 때문에 대리인의 대리권이 소멸하면 복대리인의 복대리권도 소멸한다.
\end{choices}
\end{question}
\answer{5}
